\chapter[Коррекция частотного смещения в отсчётах принятого сигнала
Переход от одной несущей к OFDM , начало формирования и приема OFDM]{Коррекция частотного смещения в отсчётах принятого сигнала
Переход от одной несущей к OFDM , начало формирования и приема OFDM}

\begin{leftbar}
    \href{https://github.com/someotb/sdr}{Ссылка на GitHub}
\end{leftbar}

\section*{Цель практики:}
Перейти от single carrier к OFDM. Сформировать OFDM символы, заполнить буффер этими символами. На приеме раскодировать принятые символы. Осуществить
алгоритм коррекции частотного смещения(CFO).

\section*{Краткие теоретические сведения}

\textbf{Переход от одной несущей к OFDM}\newline

В системах с одной несущей (Single Carrier, SC) каждый передаваемый символ
занимает всю доступную полосу частот. При прохождении через многолучевой
канал с разбросом задержек $\tau_{max}$, превышающим длительность символа,
возникает межсимвольная интерференция (ISI), устранение которой во временной
области требует сложного эквалайзера (например, адаптивного КИХ-фильтра
большой длины).

В OFDM (Orthogonal Frequency Division Multiplexing) поток данных
разбивается на $N$ параллельных потоков, каждый из которых модулирует одну
из $N$ ортогональных поднесущих. Формирование временного OFDM-символа из
блока комплексных символов данных $X[k]$, $k = 0,\dots,N-1$, выполняется
обратным дискретным преобразованием Фурье (IDFT/IFFT):

\begin{equation}
    x[n] = \frac{1}{N}\sum_{k=0}^{N-1} X[k]\, e^{j2\pi kn/N}, \quad n = 0,\dots,N-1.
\end{equation}

Для борьбы с межсимвольной интерференцией к каждому символу спереди
добавляется циклический префикс (CP) длиной $N_{CP} > \tau_{max}$ —
копия последних $N_{CP}$ отсчётов символа:

\begin{equation}
    x_{CP}[n] = x[N + n], \quad n = -N_{CP},\dots,-1.
\end{equation}

CP превращает линейную свёртку сигнала с импульсной характеристикой канала
в круговую в пределах длительности символа. Благодаря этому после удаления
CP на приёмной стороне и выполнения прямого DFT (FFT) канал в частотной
области представляется как поэлементное (плоское) умножение на комплексный
коэффициент $H[k]$ для каждой поднесущей:

\begin{equation}
    Y[k] = H[k]\,X[k] + W[k],
\end{equation}

что позволяет выполнять эквализацию однократным делением на оценку $H[k]$,
полученную по пилотным поднесущим.

\vspace{20pt}

\textbf{Коррекция частотного смещения (CFO)}\newline

Рассогласование частот гетеродинов передатчика и приёмника
(Carrier Frequency Offset, CFO) $\Delta f$ приводит к тому, что принятый
дискретный сигнал приобретает линейно нарастающий по времени фазовый
набег:

\begin{equation}
    r[n] = s[n] \cdot e^{j2\pi \Delta f \, n T_s} + w[n],
    \label{eq:cfo_model}
\end{equation}

где $T_s = 1/f_s$ — период дискретизации, $f_s$ — частота дискретизации,
$s[n]$ — переданный сигнал без искажений. Некомпенсированное CFO нарушает
ортогональность поднесущих и приводит к межканальной интерференции (ICI)
после FFT.

\vspace{20pt}

\textbf{Оценка CFO по циклическому префиксу}\newline

Поскольку циклический префикс является точной копией хвостовой части
OFDM-символа, сдвинутой на $N$ отсчётов, эти два фрагмента сигнала (в
отсутствие шума и при статичном CFO) отличаются только фазовым набегом,
накопленным за $N$ отсчётов согласно \eqref{eq:cfo_model}. Это позволяет
оценить $\Delta f$ без отдельной обучающей последовательности —
корреляцией CP с соответствующим ему фрагментом в конце символа:

\begin{equation}
    C = \sum_{k=0}^{N_{CP}-1} r^{*}[n_0 + k]\, r[n_0 + k + N],
\end{equation}

где $n_0$ — начало символа. Фаза этой корреляции связана с CFO
соотношением

\begin{equation}
    \arg(C) = 2\pi \Delta f \, N T_s,
\end{equation}

откуда оценка частотного смещения:

\begin{equation}
    \hat{\Delta f} = \frac{\arg(C)}{2\pi N T_s} = \frac{\arg(C)\, f_s}{2\pi N}.
    \label{eq:cfo_est}
\end{equation}

Для повышения точности оценки корреляционные суммы когерентно
накапливаются по всем $M$ символам в обрабатываемом блоке:

\begin{equation}
    C_{\Sigma} = \sum_{m=0}^{M-1} \sum_{k=0}^{N_{CP}-1}
        r^{*}[n_m + k]\, r[n_m + k + N],
\end{equation}

что снижает дисперсию оценки, но требует постоянства $\Delta f$ на
протяжении всего блока.

Поскольку $\arg(\cdot)$ определён на интервале $(-\pi, \pi]$, а фазовый
набег в \eqref{eq:cfo_est} накапливается за $N$ отсчётов, метод даёт
однозначную оценку только при

\begin{equation}
    |\Delta f| < \frac{f_s}{2N}.
\end{equation}

При большем смещении возникает неоднозначность (алиасинг оценки).

\vspace{20pt}

\textbf{Коррекция CFO}\newline

Компенсация выполняется умножением каждого отсчёта на комплексную
экспоненту с противоположным фазовым набегом:

\begin{equation}
    \hat{s}[n] = r[n] \cdot e^{-j2\pi \hat{\Delta f}\, n T_s}.
\end{equation}

Практически это реализуется рекурсивно (по типу цифрового генератора,
NCO), чтобы фаза коррекции непрерывно накапливалась от отсчёта к отсчёту
и не сбрасывалась на границах обрабатываемых блоков:

\begin{equation}
    \varphi[n] = \varphi[n-1] - 2\pi \hat{\Delta f} T_s,
    \qquad
    \hat{s}[n] = r[n]\, e^{j\varphi[n]},
\end{equation}

с приведением $\varphi[n]$ к интервалу $(-\pi, \pi]$ для численной
устойчивости. Такой подход эквивалентен прямому умножению на
$e^{-j2\pi\hat{\Delta f} n T_s}$, но не требует хранения абсолютного
индекса отсчёта $n$ с начала передачи.

\section*{Выполнение}

Параметры OFDM который я выбрал:
\begin{itemize}
	\item \textbf{128} поднесущих
	\item \textbf{32} циклический префик
	\item \textbf{27 + 28} защитных нулей по краям OFDM символа(в частотной области)
	\item \textbf{4} пилотные поднесущие(можно менять кол-во)
\end{itemize}

Функция которая отвечает за построение OFDM символа выглядит так:

\begin{lstlisting}[language=C++, caption={Функция, которая строит один OFDM символ}]
void build_ofdm_symbol(std::vector<uint8_t> &in, FFT_Context &context, const sharedData &sd, int &offset)
{
    context.N = sd.subcarrier;
    for (int k = 0; k < context.N; ++k)
    {
        if (sd.is_zeros[k])
        {
            context.in[k][0] = 0;
            context.in[k][1] = 0;
        }
        else if (sd.is_pilot[k])
        {
            context.in[k][0] = 1.0f / std::sqrt(2.0f);
            context.in[k][1] = 0.0f / std::sqrt(2.0f);
        }
        else
        {
            auto s = map_symbol(in, sd.modul_type_TX, offset);
            context.in[k][0] = s.real();
            context.in[k][1] = s.imag();
        }
    }

    context.ifft();
}
\end{lstlisting}

В функции используется вектора is\_pilot/is\_zeros, которые выдают true значение, если в данном индексе должен стоять защитный ноль или пилотная поднесущая.
Пилот делается как BPSK символ по стандарту 3GPP. Функция map\_symbol возвращает символ, который был сделан из входного массива бит по коду Грея, также по стандарту 3GPP.

\begin{lstlisting}[language=C++, caption=Функция которая сторит символ из битов]
std::complex<float> map_symbol(std::vector<uint8_t> &in, ModulationType mod, int &offset)
{
    if ((size_t)offset > in.size() - 8)
        return {0.0f, 0.0f};

    switch (mod)
    {
    case ModulationType::BPSK:
    {
        float v = 1.0f - 2.0f * in[offset]; offset++;
        return std::complex<float>(v, v) / std::sqrt(2.0f);
    }
    case ModulationType::QPSK:
    {
        float real = 1.0f - 2.0f * in[offset]; offset++;
        float imag = 1.0f - 2.0f * in[offset]; offset++;
        return std::complex<float>(real, imag) / std::sqrt(2.0f);
    }
    case ModulationType::QAM16:
    {
        int b0 = in[offset]; offset++;
        int b1 = in[offset]; offset++;
        int b2 = in[offset]; offset++;
        int b3 = in[offset]; offset++;
        float real = (1.0f - 2.0f * b0) * (2.0f - (1.0f - 2.0f * b2));
        float imag = (1.0f - 2.0f * b1) * (2.0f - (1.0f - 2.0f * b3));
        return std::complex<float>(real, imag) / std::sqrt(10.0f);
    }
    case ModulationType::QAM64:
    {
        int b0 = in[offset]; offset++;
        int b1 = in[offset]; offset++;
        int b2 = in[offset]; offset++;
        int b3 = in[offset]; offset++;
        int b4 = in[offset]; offset++;
        int b5 = in[offset]; offset++;
        float real = (1.0f - 2.0f * b0) * (4.0f - (1.0f - 2.0f * b2) * (2.0f - (1.0f - 2.0f * b4)));
        float imag = (1.0f - 2.0f * b1) * (4.0f - (1.0f - 2.0f * b3) * (2.0f - (1.0f - 2.0f * b5)));
        return std::complex<float>(real, imag) / std::sqrt(42.0f);
    }
    case ModulationType::QAM256:
    {
        int b0 = in[offset]; offset++;
        int b1 = in[offset]; offset++;
        int b2 = in[offset]; offset++;
        int b3 = in[offset]; offset++;
        int b4 = in[offset]; offset++;
        int b5 = in[offset]; offset++;
        int b6 = in[offset]; offset++;
        int b7 = in[offset]; offset++;

        float real = (1.0f - 2.0f * b0) * (8.0f - (1.0f - 2.0f * b2) * (4.0f - (1.0f - 2.0f * b4) * (2.0f - (1.0f - 2.0f * b6))));

        float imag = (1.0f - 2.0f * b1) * (8.0f - (1.0f - 2.0f * b3) * (4.0f - (1.0f - 2.0f * b5) * (2.0f - (1.0f - 2.0f * b7))));

        return std::complex<float>(real, imag) / std::sqrt(170.0f);
    }

    default:
        throw std::runtime_error("Unsupported mod");
    }
}
\end{lstlisting}

\vspace{20pt}

Также у меня есть функция cfo\_est, которая отвечает за коррекцию частотного смещения в пределах 7.5 кГц, так как у меня при sample rate = 1.92e6 Гц и кол-во
поднесущих равному 128, получается что каждая поднесущая расположена друг от друга на 15 кГц, отсюда и предел в 7.5 кГц.

\begin{lstlisting}[language=C++, caption=Функция коррекции частотного смещения]
std::vector<std::complex<float>> cfo_est(const std::vector<std::complex<float>> &signal, sharedData &sd)
{
    int N = sd.subcarrier;
    int CP = sd.cyclic_prefex;
    int sym_len = N + CP;
    float fs = sd.sample_rate;

    int n_symbols = signal.size() / sym_len;
    if (n_symbols < 1)
        return {};

    std::complex<float> corr_sum = 0;
    for (int sym = 0; sym < n_symbols; sym++)
    {
        int sym_start = sym * sym_len;
        std::complex<float> corr = 0;
        for (int k = 0; k < CP; k++)
            corr += conj(signal[sym_start + k]) * signal[sym_start + k + N];
        corr_sum += corr;
    }

    float epsilon = std::arg(corr_sum) / (2 * M_PI);
    float delta_f = epsilon * fs / N; // На сколько символ отклоняется от эталона по фазе в Гц

    if ((int)sd.cfo_offset.size() > sd.mtu)
    {
        sd.cfo_offset.erase(sd.cfo_offset.begin());
        sd.cfo_offset.push_back(delta_f);
    }
    else
        sd.cfo_offset.push_back(delta_f);

    std::vector<std::complex<float>> corrected(signal.size());
    for (size_t n = 0; n < signal.size(); n++)
    {
        corrected[n] = signal[n] * std::polar(1.0f, sd.current_phase);
        sd.current_phase -= 2 * M_PI * delta_f / fs;
        if (sd.current_phase < -M_PI) sd.current_phase += 2 * M_PI;
        if (sd.current_phase > M_PI)  sd.current_phase -= 2 * M_PI;
    }

    return corrected;
}
\end{lstlisting}

Функция оценивает частотное смещение по фазовому сдвигу между циклическим префиксом и его копией в конце символа, а затем разворачивает фазу каждого отсчёта на компенсирующий угол.

Сначала сигнал режется на символы длиной `N + CP` каждый. Для каждого символа берутся первые `CP` отсчётов (это циклический префикс) и сравниваются с отсчётами, которые находятся ровно на `N` позиций дальше — это как раз тот кусок символа, копией которого является CP. Сравнение делается через `conj(signal[k]) * signal[k+N]` — это комплексное произведение, которое одновременно даёт и то, насколько похожи два отсчёта по амплитуде, и на какой угол один повёрнут относительно другого по фазе. Все такие произведения внутри символа складываются, и все символы блока тоже складываются в одну сумму `corr\_sum` — это накопление увеличивает точность оценки, потому что случайный шум при суммировании частично гасится, а полезный фазовый сдвиг (если он одинаковый во всех символах) складывается когерентно.

Дальше `arg(corr\_sum)` берёт угол этого накопленного комплексного числа — это и есть средний фазовый сдвиг между CP и его оригиналом, накопленный за `N` отсчётов. Деление на `N` и умножение на `fs` переводит этот угол из "радиан за N отсчётов" в физическую величину — герцы смещения частоты, `delta\_f`.

Буфер `sd.cfo\_offset` — это просто скользящее окно последних оценок `delta\_f`, ограниченное размером `mtu`, старые значения выбрасываются с начала. В приведённом коде оно нигде не используется дальше, то есть накопленная история пока просто хранится, а не влияет на компенсацию.

Сама компенсация — это отдельный проход по всем отсчётам сигнала. Каждый отсчёт умножается на `exp(j * current\_phase)` (через `polar(1.0, phase)`, что и есть комплексное число единичной длины с этим углом). После каждого отсчёта фаза `current\_phase` уменьшается на маленький шаг `2π * delta\_f / fs` — это и есть та же логика, что "рассинхронизация растёт с каждым отсчётом", только в обратную сторону, компенсирующую эту рассинхронизацию. Условия с `-M\_PI`/`M\_PI` просто не дают углу уйти за пределы одного оборота окружности, потому что `polar` работает с любым углом одинаково, а держать угол в разумных пределах — вопрос численной аккуратности.

\section*{Вывод}
Был осуществлен переход от SC к OFDM. Написан алгоритм коррекции частотного смещения.
